\documentclass[notheorems,hyperref={bookmarks=true}]{beamer}

% --- CÁC GÓI LỆNH CẦN THIẾT ---
\usepackage[framesassubsections]{beamerprosper}
\usepackage{multirow}
\usepackage[utf8]{inputenc} 
\usepackage[utf8]{vietnam} 
\usepackage{ulem}
\usepackage{pifont}
\usepackage[mathscr]{eucal}
\usepackage{amsmath, amsthm, amssymb, amsxtra, latexsym, amscd, makeidx, tikz}
\usetikzlibrary{calc, shapes.geometric, positioning}
\usepackage{graphicx}
\usepackage{array}
\usepackage{multicol}
\usepackage{multimedia}
\usepackage{xcolor}
\usepackage{booktabs} 
\usepackage[table]{xcolor}
\usetikzlibrary{shadows}

% --- CÀI ĐẶT GIAO DIỆN ---
\usetheme{CambridgeUS}
\usecolortheme{rose} 
\usefonttheme[onlylarge]{structurebold}
\usefonttheme{professionalfonts}

\setbeamercolor{title}{fg=red!80!black,bg=blue!20!white}
\setbeamertemplate{navigation symbols}{}
\setbeamertemplate{blocks}[rounded][shadow=true]

% Logo - uncomment nếu có file
% \logo{\includegraphics[height=1cm]{logoHust.png}}

\setbeamertemplate{footline}{
  \leavevmode%
  \hbox{%
    \begin{beamercolorbox}[wd=.33\paperwidth,ht=2.5ex,dp=1.5ex,center]{author in head/foot}%
      Nguyễn Tuấn Tài
    \end{beamercolorbox}%
    \begin{beamercolorbox}[wd=.34\paperwidth,ht=2.5ex,dp=1.5ex,center]{title in head/foot}%
      \footnotesize{AI Agent hỗ trợ học ngoại ngữ}
    \end{beamercolorbox}%
    \begin{beamercolorbox}[wd=.33\paperwidth,ht=2.5ex,dp=1.5ex,right]{date in head/foot}%
      \insertframenumber/\inserttotalframenumber \hspace*{2ex} 
    \end{beamercolorbox}%
  }%
  \vskip0pt%
}

\title[AI Agent dạy ngoại ngữ]{
    \Large\bfseries Xây dựng AI Agent hỗ trợ học ngoại ngữ\\
    \large Tóm tắt kết quả chính
}

\institute[\scriptsize Khoa Toán - Tin]{
    \small
    \begin{tabular}{l}
        \textbf{Giảng viên hướng dẫn: \textcolor{red!80!black}{TS. Tạ Anh Sơn}} \\
        \textbf{Sinh viên thực hiện: \textcolor{red!80!black}{Nguyễn Tuấn Tài (20216959)}} \\
    \end{tabular} \\[0.5em]
    \textbf{Khoa Toán - Tin, Đại học Bách khoa Hà Nội}
}

\date[Hà Nội, 06/2026]{\scriptsize\bfseries\color{red!80!black}Hà Nội, 06/2026}

\begin{document}

\begin{frame}
    \maketitle
\end{frame}

\footnotesize

% ==========================================
% SLIDE 2: ĐẶT VẤN ĐỀ
% ==========================================

\begin{frame}{Đặt vấn đề: 3 Hạn chế của phương pháp truyền thống}
\begin{columns}
    \begin{column}{0.6\textwidth}
        \begin{block}{Vấn đề 1: Chi phí cao}
            Gia sư/Trung tâm: 200k-500k/giờ
        \end{block}
        
        \begin{block}{Vấn đề 2: Thiếu cá nhân hóa}
            Không phân tích lỗi chi tiết
        \end{block}
        
        \begin{block}{Vấn đề 3: Bài tập hạn chế}
            Bài tập cố định, không thích ứng
        \end{block}
        
        \vspace{0.3cm}
        
        \begin{alertblock}{✅ Giải pháp: AI Agent}
            Miễn phí, 24/7, cá nhân hóa sâu
        \end{alertblock}
    \end{column}
    \begin{column}{0.4\textwidth}
        \centering
        \begin{tikzpicture}[scale=0.7]
            \node[circle, fill=red!10, draw=red!70, thick, minimum size=2cm] (A) at (90:2.5) {\footnotesize Gia sư};
            \node[circle, fill=red!10, draw=red!70, thick, minimum size=2cm] (B) at (210:2.5) {\footnotesize Ứng dụng};
            \node[circle, fill=red!10, draw=red!70, thick, minimum size=2cm] (C) at (330:2.5) {\footnotesize Trung tâm};
            \node[circle, draw=blue!70, fill=blue!10, very thick, dashed, minimum size=2.5cm] (Goal) at (0,0) {\textbf{GIA SƯ AI}};
            \draw[->, blue!60, thick] (A) -- (Goal);
            \draw[->, blue!60, thick] (B) -- (Goal);
            \draw[->, blue!60, thick] (C) -- (Goal);
        \end{tikzpicture}
    \end{column}
\end{columns}
\end{frame}

% ==========================================
% SLIDE 3: KIẾN TRÚC MULTI-AGENT
% ==========================================

\begin{frame}{Kiến trúc Multi-Agent}
\begin{figure}[H]
\centering
\begin{tikzpicture}[
    node distance=1.5cm,
    agent_box/.style={
        rectangle,
        draw=gray!60,
        fill=blue!8,
        text width=3.5cm,
        align=center,
        minimum height=2cm,
        rounded corners=5pt,
        font=\small
    }
]

% 3 Agents
\node[agent_box] (ErrorAgent) at (0, 0) {
    \textbf{\color{red!80}Error Analyzer}\\[0.2cm]
    Phân tích lỗi
};

\node[agent_box, right=of ErrorAgent, fill=green!8] (ExerciseAgent) {
    \textbf{\color{green!80}Exercise Generator}\\[0.2cm]
    Sinh bài tự động
};

\node[agent_box, right=of ExerciseAgent, fill=purple!8] (WritingAgent) {
    \textbf{\color{purple!80}Writing Evaluator}\\[0.2cm]
    Chấm bài viết
};

% AI Tutor
\node[
    rectangle,
    draw=orange!60,
    fill=orange!15,
    text width=2.5cm,
    align=center,
    minimum height=1.5cm,
    rounded corners=5pt,
    font=\small,
    below=2cm of ExerciseAgent
] (AITutor) {
    \textbf{\color{orange!80}AI Tutor}\\
    (Điều phối)
};

% Arrows
\draw[->, >=stealth, thick, gray!70] (ErrorAgent.south) -- (AITutor.north west);
\draw[->, >=stealth, thick, gray!70] (ExerciseAgent.south) -- (AITutor.north);
\draw[->, >=stealth, thick, gray!70] (WritingAgent.south) -- (AITutor.north east);

\end{tikzpicture}
\end{figure}

\begin{alertblock}{3 Agents chuyên biệt được điều phối bởi AI Tutor trung tâm}
\end{alertblock}
\end{frame}

% ==========================================
% SLIDE 4: KẾT QUẢ 1 - HỆ THỐNG NỘI DUNG
% ==========================================

\begin{frame}{Kết quả 1: Hệ thống nội dung A1}
\begin{columns}
    \begin{column}{0.5\textwidth}
        \begin{exampleblock}{20 chủ đề A1}
            \textbf{Mỗi chủ đề = 4 bài học}
            \begin{enumerate}
                \item Grammar
                \item Vocabulary
                \item Practice
                \item Quiz
            \end{enumerate}
        \end{exampleblock}
        
        \vspace{0.3cm}
        
        \begin{alertblock}{Kết quả}
            \centering
            \Huge \textbf{80}\\
            \small bài học đầy đủ
        \end{alertblock}
    \end{column}
    \begin{column}{0.5\textwidth}
        \centering
        \begin{tikzpicture}
            \node[rectangle, draw=blue!70, fill=blue!10, rounded corners, minimum width=4cm, minimum height=1.5cm, align=center] {
                \textbf{20 topics × 4 lessons}\\
                \textbf{= 80 bài học}
            };
        \end{tikzpicture}
        
        \vspace{0.5cm}
        
        \begin{exampleblock}{Mở rộng}
            \footnotesize
            Hệ thống có khả năng mở rộng\\
            lên \textbf{190 topics} (A1-C2)\\
            với \textbf{Exercise Generator}
        \end{exampleblock}
    \end{column}
\end{columns}
\end{frame}

% ==========================================
% SLIDE 5: KẾT QUẢ 2 - ERROR ANALYZER
% ==========================================

\begin{frame}{Kết quả 2: Phân tích lỗi thông minh}
\begin{columns}
    \begin{column}{0.48\textwidth}
        \begin{exampleblock}{2 Chức năng phân tích}
            \textbf{1. Phân tích bài tập}\\
            {\footnotesize (Quiz/Practice)}
            
            \vspace{0.2cm}
            
            \textbf{2. Sửa lỗi chat}\\
            {\footnotesize (Câu văn tự do)}
        \end{exampleblock}
        
        \vspace{0.2cm}
        
        \begin{block}{Kết hợp Rule + AI}
            {\footnotesize
            ✅ Phát hiện nhanh (Rule)\\
            ✅ Phân tích sâu (AI)\\
            ✅ Độ chính xác: 92\%
            }
        \end{block}
        
        \vspace{0.2cm}
        
        \begin{alertblock}{Lợi ích}
            {\footnotesize
            ✅ Biết điểm yếu\\
            ✅ Lưu lịch sử\\
            ✅ Gợi ý luyện tập
            }
        \end{alertblock}
    \end{column}
    \begin{column}{0.52\textwidth}
        \centering
        \textbf{Ví dụ 1: Quiz}\\[0.15cm]
        \fbox{\begin{minipage}{0.95\textwidth}
            \footnotesize
            Câu hỏi: There \_\_\_ three chairs\\[0.1cm]
            User chọn: \textcolor{red}{is} (sai)\\[0.1cm]
            → AI nhận diện: "There is/are"
        \end{minipage}}
        
        \vspace{0.4cm}
        
        \textbf{Ví dụ 2: Chat}\\[0.15cm]
        \fbox{\begin{minipage}{0.95\textwidth}
            \footnotesize
            User: "I go to school yesterday"\\[0.1cm]
            Sửa: "I \textcolor{green}{went} to school..."\\[0.1cm]
            → AI nhận diện: "Past Simple"\\[0.1cm]
            → Gợi ý luyện tập thêm
        \end{minipage}}
    \end{column}
\end{columns}
\end{frame}

% ==========================================
% SLIDE 6: KẾT QUẢ 3 - EXERCISE GENERATOR
% ==========================================

\begin{frame}{Kết quả 3: Sinh bài tập tự động}
\begin{columns}
    \begin{column}{0.48\textwidth}
        \begin{exampleblock}{Exercise Generator}
            \textbf{Sinh bài tập bằng AI}\\
            {\footnotesize
            - Không giới hạn số lượng\\
            - Mỗi lần tạo mới\\
            - Thích ứng điểm yếu
            }
        \end{exampleblock}
        
        \vspace{0.2cm}
        
        \begin{block}{Ràng buộc format}
            {\footnotesize
            ✅ 6 loại bài tập\\
            ✅ Schema cố định\\
            ✅ Validation tự động
            }
        \end{block}
        
        \vspace{0.2cm}
        
        \begin{alertblock}{Ưu điểm}
            {\footnotesize
            ✅ Nội dung đa dạng\\
            ✅ Cá nhân hóa\\
            ✅ Mở rộng dễ dàng
            }
        \end{alertblock}
    \end{column}
    \begin{column}{0.52\textwidth}
        \centering
        \textbf{Ví dụ:}\\[0.15cm]
        \fbox{\begin{minipage}{0.95\textwidth}
            \footnotesize
            \textbf{User:} "Tôi muốn luyện Past Simple"\\[0.1cm]
            \textbf{AI Tutor gọi Exercise Generator:}\\
            - Chủ đề: Past Simple\\
            - Level: A2\\
            - Điểm yếu: irregular verbs\\[0.1cm]
            \textbf{AI sinh 5 bài:}\\
            1. Yesterday, I \_\_\_ (go) to school\\
            2. She \_\_\_ (eat) pizza last night\\
            3. They \_\_\_ (see) a movie...\\[0.1cm]
            → Mỗi lần tạo bài MỚI
        \end{minipage}}
    \end{column}
\end{columns}
\end{frame}

% ==========================================
% SLIDE 7: KẾT QUẢ 4 - WRITING EVALUATOR
% ==========================================

\begin{frame}{Kết quả 4: Đánh giá bài viết tự động}
\begin{columns}
    \begin{column}{0.45\textwidth}
        \begin{exampleblock}{Writing Evaluator}
            \textbf{Chấm 4 tiêu chí CEFR:}
            \begin{enumerate}
                \item Ngữ pháp
                \item Từ vựng
                \item Nội dung
                \item Cấu trúc
            \end{enumerate}
        \end{exampleblock}
        
        \vspace{0.3cm}
        
        \begin{alertblock}{Phản hồi}
            ✅ Chỉ từng câu sai\\
            ✅ Gợi ý cải thiện\\
            ✅ Lưu lịch sử
        \end{alertblock}
    \end{column}
    \begin{column}{0.55\textwidth}
        \centering
        \textbf{Ví dụ kết quả:}\\[0.2cm]
        \fbox{\begin{minipage}{0.9\textwidth}
            \footnotesize
            \textbf{Topic:} My Daily Routine (A1)\\[0.2cm]
            \begin{tabular}{lr}
                \textbf{Grammar:} & 7/10 \\
                \textbf{Vocabulary:} & 8/10 \\
                \textbf{Content:} & 9/10 \\
                \textbf{Structure:} & 7/10 \\
                \hline
                \textbf{Tổng:} & \textbf{7.75/10}
            \end{tabular}
            \\[0.3cm]
            \textbf{Nhận xét:} Bài viết tốt,\\cần chú ý thì động từ câu 3, 5.
        \end{minipage}}
    \end{column}
\end{columns}
\end{frame}

% ==========================================
% SLIDE 8: KẾT QUẢ 5 - BỘ NHỚ + TRIỂN KHAI
% ==========================================

\begin{frame}{Kết quả 5: Bộ nhớ học tập \& Triển khai}
\begin{columns}
    \begin{column}{0.5\textwidth}
        \begin{block}{Bộ nhớ 2 tầng}
            \textbf{Short-term:}\\
            {\footnotesize Duy trì hội thoại}
            
            \vspace{0.2cm}
            
            \textbf{Long-term:}\\
            {\footnotesize Lưu lỗi sai lâu dài (PostgreSQL)}
        \end{block}
        
        \vspace{0.3cm}
        
        \begin{alertblock}{→ Cá nhân hóa sâu}
            AI "nhớ" điểm yếu từng user
        \end{alertblock}
    \end{column}
    \begin{column}{0.5\textwidth}
        \begin{block}{Triển khai Production}
            \begin{itemize}
                \item ✅ FastAPI + LangChain
                \item ✅ Streamlit
                \item ✅ PostgreSQL (Neon Tech)
                \item ✅ Groq API (Llama 3.1)
                \item ✅ Render.com
            \end{itemize}
        \end{block}
        
        \vspace{0.2cm}
        
        \begin{exampleblock}{Live}
            \tiny \texttt{lang-prj.onrender.com}
        \end{exampleblock}
    \end{column}
\end{columns}
\end{frame}

% ==========================================
% SLIDE 9: TỔNG KẾT
% ==========================================

\begin{frame}{Tổng kết}
\begin{columns}
    \begin{column}{0.5\textwidth}
        \begin{exampleblock}{Giải quyết 3 vấn đề}
            \begin{itemize}
                \item ✅ Chi phí/thời gian
                \item ✅ Cá nhân hóa
                \item ✅ Hỗ trợ linh hoạt
            \end{itemize}
        \end{exampleblock}
        
        \vspace{0.3cm}
        
        \begin{block}{5 Kết quả kỹ thuật}
            \footnotesize
            \begin{enumerate}
                \item Multi-Agent
                \item 950 bài học
                \item Error Analyzer
                \item Writing Evaluator
                \item Production Cloud
            \end{enumerate}
        \end{block}
    \end{column}
    \begin{column}{0.5\textwidth}
        \begin{alertblock}{Hạn chế}
            Chỉ Đọc/Viết ở cơ bản
        \end{alertblock}
        
        \vspace{0.3cm}
        
        \begin{exampleblock}{Hướng phát triển}
            \begin{itemize}
                \item Voice AI (Nghe/Nói)
                \item Mở rộng 4 kỹ năng CEFR
                \item Scoring model riêng
            \end{itemize}
        \end{exampleblock}
    \end{column}
\end{columns}
\end{frame}

% ==========================================
% SLIDE 10: CẢM ƠN
% ==========================================

\begin{frame}
    \begin{center}
         \Huge\bfseries Cảm ơn thầy cô\\và các bạn!
         
         \vspace{1cm}
         
         {\Large \textcolor{blue!80}{AI Agent hỗ trợ học ngoại ngữ}}\\[0.5cm]
         {\large Nguyễn Tuấn Tài - 20216959}
    \end{center}
\end{frame}

\end{document}
